
\chapter{Free diffusion}

\section{Formulation}

Let us consider a pure diffusion equation void of any drift term, also called the heat equation.

\begin{align*}
\partial_t u - D \Delta u &= 0 \\
u(x) &= u_0(x) \\
u(x,t) &= 0
\end{align*}
where
\begin{align*}
\Omega &= [0,1]^d\\
T&=5.0 \\ 
D &= \frac{\bar{D}}{d}, \quad \bar{D}=0.01 \\ 
u_0(x) &=
\prod_{i=1}^d \sin(\pi x_i)
+
\frac{0.67}{\sqrt{d}}
\sum_{i=1}^d \sin(2\pi x_i) \left(\prod_{k=1,k\neq i}^d \sin(\pi x_k)\right).\\
\end{align*}


The corresponding analytical solution is then given by
%(see the Appendix for the derivation)

$$
u(t,x) =
e^{-\gamma t}
\prod_{i=1}^d \sin(\pi x_i)
+
e^{-\delta t}
\frac{0.67}{\sqrt{d}}
\sum_{i=1}^d \sin(2\pi x_i) \left(\prod_{k=1,k\neq i}^d \sin(\pi x_k)\right)
$$
with
\begin{itemize}
\item $\gamma = D d \pi^2 = \bar{D} \pi^2$
\item $\delta = D (d+3) \pi^2 = \bar{D} (1+\frac{3}{d}) \pi^2\:$.
\end{itemize}


Apart from the simple product of isotropic sines,
the solution also contains a sum of higher-frequency terms,
which slightly increases its spatial and temporal complexity.


\section{Behavior in general dimension}

In general dimension, the solution has a peak in the positive quadrant near the origin,
which later shifts towards the center of the hypercube, since $\gamma$ decays more slowly than $\delta$.

%$$
%u((0.5,\dots,0.5), t) = e^{-\gamma t} = e^{-\bar{D} \pi^2}
%$$
\textbf{Decay rate $\gamma$}

The decay rate $\gamma$ of the leading term is independent of dimension.
If $D$ were just a constant independent of $d$, the solution would decay exponentially with dimension,
leading to negligible values even for relatively modest values of $d$.
%Hence, we consider $D=\frac{\bar{D}}{d}$.

\textbf{$L^2$ norm}
$$
||u(t,\cdot)||_{L^2(\Omega)}^2 =
\frac{1}{2^d}
\left(
e^{-\gamma t}
+
0.67^2
e^{-\delta t}
\right)
$$

We see that the $L^2$ norm decreases exponentially with dimension.
If we were to consider the $L^2$ error as a measure of accuracy,
then a trivial prediction $u=0$ would appear better and better with higher dimensions.
%then for a trivial case where the predicted solution were 0,
%this would make the trivial solution appear artificially good in higher dimensions.
Instead, we shall prefer the relative $L^2$ error as a metric, see below.

%Note, that the above exponential dependend on $d$
%could be modified by multiplying the initial condition by $2^{d/2}$, but that would
%inevitably lead to an exponential concentration of mass, so we keep the norm as it is.




%- the notorious peak of the orange has more volume that its core (see Appendix)

\section{Metrics in higher dimensions}

As dimensionality $d$ increases, the geometry of the space changes drastically,
and the volume shifts aggressively toward the boundaries.


\textbf{A. Relative $L^2$ error}
%IC,TC pred vs true, if analytical known: PDE pred vs true
$$
\frac{
\sum_{i=1}^n(\hat y_i - y_i)^2
}{
\sum_{i=1}^n(y_i)^2
}
$$
where $\hat y_i$ are the predictions and $y_i$ is the ground truth value.
use LHS

\textbf{B. $L^\infty$ error}
%IC,TC pred vs true, if analytic known: PDE pred vs true
$$
E_\infty = \max_\Omega\|\hat{y}_i - y_i\|
$$
%need to look at the section of the domain where the solution is to have the max value
%sample the center via gaussian?

%we report total $L^\infty$ error, we mean the max
%$$
%max(
%L^\infty_{pde},\,
%L^\infty_{bc},\,
%L^\infty_{ic}
%)
%$$
%by relative $L^2$, we mean the average
%$$
%\frac{
%L^2{pde}+
%L^2{bc}+
%L^2{ic}
%}{3}
%$$
%in this construction the 'relative' keeps its meaning.
%Since comparing a zero prediciton with a nonzero solution gives 1.

Together, $L^\infty$ tells us about point error, rel. $L^2$ more about domain error.


In the experiments below,
to enter the curse of dimensionality regime more quickly,
we proceed by 2 dimensions at a time.

In addition, for both methods, our goal will be to reach the $L^\infty$ and relative $L^2$ errors of order $1e-3$.


\section{Sparse grid combination technique}

As a classical baseline, we consider the sparse grid combination technique.
For the time discretization, we use the Crank--Nicolson scheme.
On each component grid, this yields a sparse linear system, which we solve by GMRES without preconditioning.
We also tested CG with ILU preconditioning, but in our experiments GMRES consistently reached the same accuracy faster,
and will therefore be employed throughout this section.
Since the analytical solution is available, we evaluate the approximation error directly at randomly sampled points using Latin hypercube sampling,
and report both the relative $L^2$ error and the $L^\infty$ error.
Due to the parallel nature of the algorithm, the sampled points cannot be scattered uniformly at random
in the temporal domain, as that would require frequent interruption of the parallel solvers, thus significantly
hindering the main advantage of the method. Instead, we select a set of time slices $\{1.25, 2.5, 3.75, 5\}$
and compute the error metrics cumulatively.
The error at the initial condition is computed separately.


\subsection{2-dimensional calibration case}

We begin with the two-dimensional setting, where the method can be calibrated against the analytical solution at low computational cost.
We set $\Delta t=0.01$ and benchmark several sparse-grid levels $n$.
We also tested a smaller time step, $\Delta t=0.001$, but observed only a negligible change in the reported errors,
which is consistent with the good stability property of the implicit Crank--Nicolson scheme.

All runs in this subsection use 8 parallel workers.
In the 2-dimensional setting, even at larger levels, the number of component grids is still small,
and both runtime and memory footprint remain nearly constant across the tested values of $n$:
the runs finish in about $0.2$\,s and require about $100$\,MB of RAM.
The approximation accuracy provided by increasing levels is not therefore
hindered by computational cost.

Below, we report the levels $n$ that achieved either relative $L^2$ error or $L^\infty$ error of approximately $10^{-2}$ or smaller.

\begin{table}[ht]
\centering
\begin{tabular}{ |c|c|c|c|c|c|c| } 
\hline
level $n$ & rel $L^2$ & $L^\infty$ & \# grids & \# grid points \\
%\hline
%4 & 0.0241626 & 0.0506153 & 7 & 113 \\
\hline
5 & 0.00725 & 0.01559 & 9 & 257 \\
\hline
6 & 0.00210 & 0.00465 & 11 & 577 \\
\hline
7 & 0.00060 & 0.00131 & 13 & 1,281 \\
\hline
8 & 0.00017 & 0.00036 & 15 & 2,817 \\
\hline
\end{tabular}
\caption{Accuracy of the sparse grid combination technique in the two-dimensional case.}
\label{tab:sg-2d}
\end{table}

%The pattern is clear: increasing the sparse-grid level leads to a rapid reduction in both error metrics.
To our initial surprise, the largest error was achieved at the initial condition, and then decreased as $t \to T$.
However, due to the smoothing effect of the heat equation, this behavior is to be expected.

As mentioned in the chapter on the sparse-grid method,
the theoretical estimate for the approximation error is given as
$$
\|u-\widehat u_n\|_{L^\infty(\Omega)}
=
O(h_n^2\:n^{d-1}).
$$
In fact, our experiments support this estimate quite well in the asymptotic regime.
In particular, the observed decay of the $L^\infty$ error for larger values of $n$ follows the theoretical prediction closely, see Figure~\ref{fig:sg_linf_matches_theor_err}.


%[show plot of the solution at lvl=1,lvl=2,lvl=3 side by side - to see the basis functions / grid points at play]
\begin{figure}
    \centering
    % see combination technique repo: gs_heat_eq__2026-07-13--01-28-36__d=4_gmres_levels=2:8
    \includegraphics[width=\linewidth]{figures/sg_2d_basis_funs_visible.png}
    \caption{Clearly visible basis functions in the interpolation of the initial condition for levels $n=1,2,3$.}
    \label{fig:sg_2d_basis_funs_visible}
\end{figure}


\begin{figure}
    \centering
    % see combination technique repo: gs_heat_eq__2026-07-13--01-28-36__d=4_gmres_levels=2:8
    \includegraphics[width=\linewidth]{figures/sg_linf_matches_theor_err.png}
    \caption{Asymptotic behavior of the $L^\infty$ error in our experimental setup compared with the theoretical result ($d=4$).}
    \label{fig:sg_linf_matches_theor_err}
\end{figure}



\subsection{4 dimensions}
In four dimensions, the sparse-grid approximation remains computationally manageable on modest hardware,
but the growth in computational complexity is already visible.
The number of component grids and the total number of grid points increase substantially with the level, and runtime rises from a few seconds to about a minute.
At the same time, the accuracy degrades relative to the two-dimensional case, and an error below $10^{-2}$ is only first achieved at $n=7$.

For this experiment, we increased the test set 10-fold, to $10^5$ sampled space-time points.
The results are summarized below.
\begin{table}[ht]
\centering
\begin{tabular}{ |c|c|c|c|c|c|c| } 
\hline
level $n$ & rel $L^2$ & $L^\infty$ & \# grids & \# grid points & peak RAM & runtime \\
%\hline
%4 & 0.0776054 & 0.152767 & 35 & 2769 & 113.605 & 1.30406 \\
%\hline
%5 & 0.0301881 & 0.0628298 & 69 & 7681 & 121.98 & 2.56963 \\
\hline
6 & 0.01108 & 0.02388 & 121 & 20,481 & 151 MB & 5.96 s\\
\hline
7 & 0.00389 & 0.00863 & 195 & 52,993 & 242 MB & 18.45 s\\
\hline
8 & 0.00132 & 0.00298 & 295 & 133,889 & 528 MB & 71.14 s \\
\hline
\end{tabular}
\caption{Accuracy and computational cost of the sparse grid combination technique in 4 dimensions.}
\label{tab:sg-4d}
\end{table}
Already in this case, the cost grows quite rapidly with the level.
From $n=6$ to $n=8$, the number of component grids increases from $121$ to $295$,
the peak memory footprint more than triples,
and runtime increases by more than an order of magnitude.
This behavior becomes even more pronounced in higher dimensions.


\subsection{6 dimensions}
The 6-dimensional case marks the point where the practical limitations of the combination technique become significant.
Although a few hundred component grids are still manageable, reaching acceptable accuracy requires moving to higher levels,
where both runtime and memory consumption grow rapidly.

The computations still use 8 workers, but the runtime is now on the order of minutes rather than seconds.
The obtained results are shown in Table~\ref{tab:sg-6d}.
Note that the test set was again increased 10-fold, to $10^6$ sampled space-time points.

\begin{table}[ht]
\centering
\begin{tabular}{ |c|c|c|c|c|c|c| } 
\hline
level $n$ & rel $L^2$ & $L^\infty$ & \# grids & \# grid points & peak RAM & runtime \\
%\hline
%2 & 0.537902 & 0.762118 & 7 & 3645 & 220.648 & 6.66377 \\
%\hline
%3 & 0.30475 & 0.472944 & 28 & 14337 & 250.809 & 25.8216 \\
%\hline
%4 & 0.154217 & 0.258022 & 84 & 49761 & 361.441 & 85.0774 \\
\hline
5 & 0.0668 & 0.1276 & 210 & 159,489 & 745 MB & 4 min 43 s \\
\hline
6 & 0.0287 & 0.0579 & 462 & 483,201 & 2.19 GB & 18 min 31s \\
\hline
\end{tabular}
\caption{Accuracy and computational cost of the sparse grid combination technique in six dimensions.}
    \label{tab:sg-6d-estimate}
\end{table}

Here, the cost of each level increment is substantial.
Going from $d=4$ to $d=6$ on level $n=6$ increases the number of component grids from $121$ to $462$,
the number of grid points from $20481$ to $483201$,
and the runtime from a few seconds to more than $18$ minutes.
At the same time, the peak memory usage grows beyond $2$\,GB.

To push the method to higher levels, we run the next experiments on a cluster,
this enables us to increase the number of workers from 8 to 48,
and the available RAM from 16\,GB to 72\,GB.
%rather than on a amd laptop as was done up to this point.
In addition, to achieve higher efficiency, we
avoid forming the full sparse matrix explicitly and instead represent the discretized operator via SciPy's \texttt{LinearOperator} class,
which significantly reduced the memory footprint.

The reported RAM measurements in this part should be taken only as rough indicative values, as the logging on the cluster side here was not fully reliable.

\begin{table}[ht]
\centering
\begin{tabular}{ |c|c|c|c|c|c|c| } 
\hline
level $n$ & rel $L^2$ & $L^\infty$ & \# grids & \# grid points & peak RAM & runtime \\
\hline
6 & 0.0287 & 0.0579 & 462 & 5e05 & 350 MB & 6 min \\
\hline
7 & 0.0117 & 0.0245 & 923 & 1.4e06 & 600 MB & 10 min \\ % 32 workers
\hline
8 & 0.0046 & 0.0098 & 1,709 & 4e06 & 1.5 GB & 25 min \\ % 48
\hline
9 & 0.0017 & 0.0037 & 2,975 & 1e07 & 4 GB & 45 min \\
\hline
\end{tabular}
\caption{Accuracy and computational cost of the sparse grid combination technique in six dimensions.}
\label{tab:sg-6d}
\end{table}

We see that levels 8 and 9 offer a good enough accuracy.



\subsection{8 dimensions}

Here, the growth becomes already rather extreme.
The resulting number of component grids and total grid points can no longer be handled comfortably in memory,
and on the cluster we observed out-of-memory failures even with 256\,GB of RAM on 16 workers.

To perform the computation, each worker has to be able to hold its compoent grid in memory.
And, on a fixed RAM budget, spawning many workers in parallel leads to large RAM consumption.
Hence, in order to obtain some results in this space,
one must either dramatically increase the available RAM,
or somewhat contrary to the parallel advantagecombination technique offers,
reduce the number of workers to fit the growing number of component grids and grid points in RAM.

As the cluster prohibits us from reserve more than 256 GB of RAM, the combination technnique
here becomes computationally prohibitive.


%%%%%%%%%%%%%%%%%%%%%%%%%%%%%%%%%%%%%%%%%%%%%%%%%%%%%%%%%%%%%%%%%%%%%%%%%%%%%%%%%%%%%%%%%%%%%%%%%%%%%%%%
%%%%%%%%%%%%%%%%%%%%%%%%%%%%%%%%%%%%%%%%%%%%%%%%%%%%%%%%%%%%%%%%%%%%%%%%%%%%%%%%%%%%%%%%%%%%%%%%%%%%%%%%
%%%%%%%%%%%%%%%%%%%%%%%%%%%%%%%%%%%%%%%%%%%%%%%%%%%%%%%%%%%%%%%%%%%%%%%%%%%%%%%%%%%%%%%%%%%%%%%%%%%%%%%%


\section{PINNs}

Next, we turn to physics-informed neural networks.

%When trainig neural network we frequency have to decide a the values of multiple parameters 
%In the sections below, we shall frequency use the term \emph{grid search}.
%This concept is frequently used in the machine learning community
%and it refers to theevaluating all combinations of the selected hyperparameter values.

As in the sparse-grid experiments, some of the lower-dimensional runs are carried out on a local CPU machine,
while the heavier computations are moved to the cluster and run on GPU hardware.

Every $k$-th step during training, we compute a)
the $L^\infty_{pde}$ and relative $L^2_{pde}$ errors based on a fixed, pre-sampled
set of points obtained by LHS over the whole space-time domain, and b) the $L^\infty_{ic}$ and relative $L^2_{ic}$ errors
of the initial condition based on a similar, albeit slightly smaller, fixed dataset.
Total $L^\infty$ error is computed as $\max(L^\infty_{pde}, L^\infty_{ic})$,
while the relative $L^2$ error is the average of the PDE and IC relative $L^2$ errors.

In addition, due to the significant oscillations in network performance,
and hence in the loss and error metrics computed during training,
instead of reporting the errors obtained at the very last training step, we
average the error metrics over the last 100 collected values and report that number.
This significantly reduces the effect of random oscillation and makes comparisons between runs more stable.

In addition, we will frequently use the compact notation '10k' to mean 10,000 and '1M' to mean 1,000,000.

In all the experiments below, we use the following setup unless stated otherwise:

\begin{itemize}
\item Latin hypercube sampling (LHS)
\item Learning rate (lr) of $10^{-3}$ with an exponential decay factor of $\gamma=0.9$
\item $tanh$ activation functions
\item Hard constraints for the boundary condition:
$$
u_\theta(t,x) = \hat{u}_\theta(t,x) \cdot \prod_{i=1}^d \sin(\pi x_i)
$$
\end{itemize}

Also, during training, we periodically sample the dataset so that in each step the network receives a never-before-seen set
of points. In addition, in accordance with the PINN literature, the batch for the boundary condiiton is set to 1/8 the size of the
batch for the PDE.\cite{experts_guide}


\subsection{2-dimensional calibration case}

Guided by the base PINN architecture introduced in the PINN chapter, we begin with the following hyperparameters:
\begin{center}
\begin{tabular}{ |c|c|c|c|c } 
\hline
NN size & steps & batch size (bs) & number of test points & lr decay frequency\\
\hline
4x64  & 10k & 2048 & 10k & 10\\
\hline
\end{tabular}
\end{center}
Here, `4x64` denotes a neural network with 4 hidden layers, each consisting of 64 neurons.

We also use exponential learning-rate decay, applied 10 times during training.




\subsubsection{Loss weights: $\lambda_{pde}, \lambda_{ic}$}
As discussed in the PINN chapter, the loss landscape is particularly sensitive
to the weights assigned to the individual loss terms, in our case $\lambda_{pde}$ and $\lambda_{ic}$.
To determine a reasonable choice, we first perform
a parameter sweep over several values. In particular, we fix $\lambda_{pde}=1.0$ and vary $\lambda_{ic}$,
spanning 5 orders of magnitude.

Since we expect this ration to be important, we repeat the experiment for 3 random seeds and report the averaged results.

As seen in Figure~\ref{fig:lambda_ic__GS}, the value $\lambda_{ic}=0.1$ gives the best results.
Hence we set this value for the following runs. We are not conrned with finetuning this value further, since the
optimal values is dependent on the other hyperparameter choices which we have not yet performed, and such a holistic
optimizaiton is computationally very demanding.
Thus we are content with finding a reasonable values of $\lambda_{ic}$ and set this as our base network.

\begin{figure}
    \centering
    \includegraphics[width=\linewidth]{figures/lambda_ic__GS.png}
    \caption{Influence of $\lambda_{ic}$ on the final loss and $L^\infty$ error.}
    \label{fig:lambda_ic__GS}
\end{figure}


\subsubsection{Loss landscape}
%bs: 1024 2048 4096

\begin{figure}
    \centering
    \IfFileExists{figures/pinn_2d_oscills_losses.png}{
        \includegraphics[width=\linewidth]{figures/pinn_2d_oscills_losses.png}
    }{
        \fbox{\parbox[c][0.3\textheight][c]{0.95\linewidth}{\centering Missing figure: \texttt{figures/pinn\_2d\_oscills\_losses.png}}}
    }
    \medskip
    \IfFileExists{figures/pinn_2d_oscills_rel_l2.png}{
        \includegraphics[width=\linewidth]{figures/pinn_2d_oscills_rel_l2.png}
    }{
        \fbox{\parbox[c][0.3\textheight][c]{0.95\linewidth}{\centering Missing figure: \texttt{figures/pinn\_2d\_oscills\_rel\_l2.png}}}
    }
    \caption{Oscillating behavior of the base network with batch size $=2048$ and learning-rate decay frequency $=10\times$.}
    \label{fig:pinn_2d_oscills}
\end{figure}

After plotting the resulting metrics obtained from the preceding run, see Figure~\ref{fig:pinn_2d_oscills},
we observe both the loss and, more importantly, the error metrics oscillate quite strongly.
In our experiments, this behavior appears largely irrespective of the choice of $\lambda_{ic}$.
Since the batch contains only 2048 points sampled over the full space-time domain,
we speculate whether the batch size itself is contributing to this instability.
To test this, we compare the base run with corresponding runs using batch sizes 1024 and 4096.

To our surprise,
the errors and the oscillation patterns are not strongly correlated with the batch size,
which is somewhat surprising given that the gradient updates are usually quite sensitive
to the noise level in the batches.\cite{goodfellow}

What we do notice, however, is that smaller batch sizes run much faster on CPU,
which is not too surprising given the more sequential character of the hardware.
On GPU, the situation is different. Since both the model and the dataset are still relatively small,
changing the batch size has only a minor effect on runtime on this parallel hardware.

%\begin{center}
%\begin{tabular}{ |c|c|c|c|c| } 
%\hline
%batch size & time [cpu run] & time[gpu run] & rel $L^2$ & $L^\infty$ \\
%\hline
%1024 & 108s & ?s & 0.617e-3 & 1.27e-3 \\
%\hline 
%2048 & 168s & ?s & 0.684e-3 & 1.44e-3 \\
%\hline 
%4096 & 292s & ?s & 0.7e-3 & 1.3e-3 \\
%\hline
%\end{tabular}
%\end{center}

Since some of our experiments are performed on CPU only,
we select a batch size of 1024 as the most efficient setting for the subsequent experiments.

We also observe that further decreases in batch size tend to slightly hinder the training performance.

%
%and the difference between the resampling and fixed datset regimes.
%Most efficient bs?
%
%average over 3 seeds again
%
%bs effect on training dynamics, hence apart from the error,
%we also report time.
%bs: 512 vs 2048 vs 8096
%
%We keep the the number of steps fixed (10k) - fixed n of optimizer steps.
%- bs = 512, ...
%
%%To this effect we consider two modes of grid search.
%%Second, fix the n of points that pass through the network during training.
%%- bs = 512 --> nsteps = 40k
%%- bs = 2048 --> nsteps = 10k
%%- bs = 8096 --> nsteps = 2.5k
%The batch size is the botthle neck?
%If we look at the memoty log, the GPU VRAM usage a peaked at 60MB. During the run...
%However, we see that smaller batch sizes lead to smaller $L^2$ and $L^\infty$ errors.
%The GPU needs much large batch size
%%\begin{figure}
%%    \centering
%%    \includegraphics[width=150mm]{figures/training_loss.png}
%%    \caption{loss minibatch, fixed}
%%    \label{fig:rar-d}
%%\end{figure}
%on gpu:
%even thought the number of collocation points is the same, 
%too small peak gpu mem utilization
%larger bs ?
%loss landscape ?
%does not matter that much, 1024 seem like a good compromise




We next try to further mitigate the oscillatory loss and error landscape.
The simplest way to do so is to adjust the learning-rate schedule and observe the resulting effect on both the $L^2$ and $L^\infty$ errors.

In the base architecture, the learning rate is decreased every 1000 steps, that is, 10 times during training.
We now increase this frequency to 20. We observe that the learning-rate schedule does indeed affect
the magnitude and distribution of the oscillations.

As expected, this noticeably changes the loss dynamics and reduces the infrequent high spikes.

To improve the final accuracy further, we next
increase the number of steps to 20,000 and retune the decay frequency for this longer run.
The longer training horizon then yields more than two times better results compared to the same network trained for 10k steps.

We report a good middle ground at 25, meaning that the network ends the training
at a learning rate of approximately 7\% of its starting value.


All these runs still finish within a a few minutes, so this tuning remains relatively inexpensive in 2 dimensions.

The obtained loss and error metrics are given below.

\begin{center}
\begin{tabular}{ |c|c|c|c|c|c| } 
\hline
batch size & steps & lr decay frequency & $\ell_{\mathrm{total}}$ & rel $L^2$ & $L^\infty$ \\
\hline
2048 & 10k & 10x & 1.898e-06 & 0.00344 & 0.00572 \\
\hline
1024 & 10k & 25x & 0.635e-06 & 0.00243 & 0.00411 \\
\hline
1024 & 20k & 25x & 0.125e-06 & 0.00102 & 0.00174 \\
\hline
\end{tabular}
\end{center}
% a_heat_2d_oscills/


% /home/antos_j/diploma_sketchbook/PINN/case0_HeatEq/gridsearch__2026-07-09--11-27-08_resampl_decay750_20k/seed=42__prevent_resampling=False__bs=1024__n_res_points=102400/report.json


% if we try gradnorm, the lambda is optimal - goes to 10 vs 1 but same res


%\subsubsection{LBFGS}
%
%\paragraph{LBFGS alone}
%We first consider a pure LBFGS run.
%LBFGS expects large collation size, hence we slighlty increase the
%batches of sizes to 2k and 4k. And, since one lbfgs step can consist
%of multiple line search steps, we select a relatively modes 500 steps.
%
%Super quick convergence.
%line search 0 only a few steps ater the initia period.
%After that the line search only did not find any better minimum
%
%Suprisingly, we observed no noticable change between the two batch sizes.
%What we noticed, though was a significant worsening of the
%relative $L^2$ error for the initial conditon.
%Guide by this we increased the $\lambda_{ic}$ to the default 1.0, and the 
%performaced imporved by almost an ordr of magnitude - even below the pde one, without affecting the value of pde erros.
%
%\begin{center}
%\begin{tabular}{ |c|c|c|c|c|c| } 
%\hline
%$\ell_{pde}$ & $\ell_{ic}$ & $rel L^2$ pde & $rel L^2$ ic & $L^\infty$ pde & $L^\infty$ ic \\
%\hline
%2.8e-7 & 2.77e-7 & 0.000604 & 0.0009 & 0.00134 & 0.0014 \\
%\hline
%\end{tabular}
%\end{center}
%% /home/antos_j/diploma_sketchbook/PINN/case0_HeatEq/run_latest_l3_bs4000_dec50/report.json
%
%\paragraph{Adam + LBFGS}
%We run adam for until adam 1e-5 and 1e-6.
%At end of each run, we obtain the following errors:
%
%\begin{center}
%\begin{tabular}{ |c|c|c|c|c|c| } 
%\hline
%terminal loss value & steps & $rel L^2$ pde & $rel L^2$ ic & $L^\infty$ pde & $L^\infty$ ic \\
%\hline
%$1e-5$ & 2384 & 0.00811 & 0.00924 & 0.01693 & 0.01719 \\
%\hline
%$1e-6$ & 6121 & 0.00266 & 0.00379 & 0.00563 & 0.00567 \\
%\hline
%\end{tabular}
%\end{center}
%% /home/antos_j/diploma_sketchbook/PINN/case0_HeatEq/run_latest__adam1e-5/report.json
%[show plot: adam_e-5,6_starting_points - add lines to 1e-5,6]
%
%% consider $\lambda_{ic}=0.1, 1$ or only 1.0 ?
%
%For 
%We then fix the batch size of 2048, and 4096.
%Then run lbfgs for 250 (since 500 super flat) steps with a few random seeds,
%starting either from adam 1e-5 or 1e-6.
%
%All four resulting runs give almost equivalent results.
%Regardless of whethet we start from adam 1e5 or 6 and regardless of the batch size.
%
%Almost 10x decreate in all error metrics:
%
%The most efficient choice:
%2048, 1e-5, (lambda 1.0)
%\begin{center}
%\begin{tabular}{ |c|c|c|c| } 
%\hline
%$rel L^2$ pde & $rel L^2$ ic & $L^\infty$ pde & $L^\infty$ ic \\
%\hline
%0.000503 & 0.000490 & 0.000995 & 0.000843 \\
%\hline
%\end{tabular}
%\end{center}
%% /home/antos_j/diploma_sketchbook/PINN/case0_HeatEq/gridsearch__2026-07-09--13-08-44__lbfgs_GS_adam1e-5,6_42/lambda_ic=1p0__n_res_points=204800__bs=2048__starting_model=run_latest__adam1e-5-model.pth/report.json
%
%
%
%We naturally pick the combiantion of adam for till 1e-5, followed by LBFGS for 250 steps.
%
%[figure: pure adam, pure lbfgs, mix]
%
%
%bottom line:
%- minimizer of L2
%- prefers lambda 1.0



\subsection{4 dimensions}
%Here, we reuse the same optimal net architecutre and hyperparameter choices
%decay 25
%n steps, NN, bs
%We run adam for 20k, 40k and nn 4x64 and 4x128
%testing point to 1M of the spatial temporal domain.
%In addition we try lbfgs, also for 4x64 and 4x128.
%Adam till 1e-5, followed by LBFGS for 250 steps.

%/storage/praha1/home/simulant_antos/diploma_sketchbook/PINN/case0_HeatEq/gridsearch__26-07-09--20:01:51__GS_d4/

We increase the number of space-time test points to $10^5$.

Here, we reuse the best hyperparameters found in the previous case, namely
$\lambda_{ic}=0.1$, batch size of 1024, and learning-rate decay frequency of 25.
To address the more complex space, we try to increase the network size, from 4x64 to 4x128,
effectively quadrupling its capacity. To account for this larger network,
we also consider extending the training horizon from 20k to 40k steps.

The results are shown in Figure~\ref{fig:models_d4}. Of the four tested combinations,
it is apparent that the longer training horizon leads to better results
irrespective of the model capacity. %and leads to approximately 2x lower errors.
As for the model capacity, it was somewhat surprising to see the smaller model
perform nearly as well as the larger one.
The larger network still benefits from the longer run, which suggests that increasing model size typically requires
a correspondingly longer training horizon.
We also observe that the loss correlates with the final error reasonably well, which is useful when comparing runs during training.
Below we report the two best runs, both corresponding to 40k training steps.
We also report the peak allocated GPU memory (VRAM). In both cases, the runtime is about 540\,s.
Compared with the 2-dimensional case, the error remains only moderately larger, which is a good outcome given the increased dimensionality.

For the two best-performing networks (40k):
\begin{center}
\begin{tabular}{ |c|c|c|c|c|c| } 
\hline
network size & $\ell_{\mathrm{total}}$ & rel $L^2$ & $L^\infty$ & peak allocated VRAM & runtime\\
\hline
4x128 & 3.88e-8 & 0.00162 & 0.00255 & 120 MB & 9 min\\
\hline
4x64 & 1.126e-7 & 0.00189 & 0.00304 & 71 MB & 9 min\\
\hline
\end{tabular}
\end{center}
% cluster_4d/

\begin{figure}
    \centering
    \IfFileExists{figures/nn_scaling_comparison_d=4.png}{
        \includegraphics[width=\linewidth]{figures/nn_scaling_comparison_d=4.png}
    }{
        \fbox{\parbox[c][0.3\textheight][c]{0.95\linewidth}{\centering Missing figure: \texttt{figures/nn\_scaling\_comparison\_d=4.png}}}
    }
    \caption{Error and loss for various network sizes and training durations.}
    \label{fig:models_d4}
\end{figure}

Compared to our initial expectation, the GPU memory footprint remains quite low for both models.


\subsection{6 dimensions}
% exaclty same as in d=4, but also
%causal weights and causal sampling
%RAR
%SDGD

Here, we proceed similarly as in $d=4$.
First, we increase the number of space-time test points to 1M.

We then probe the effect of increasing both the network size and the number of training steps.

We observe a similar qualitative pattern as in four dimensions.
Both increasing the model size and extending the training duration improve the final error,
but the effect of longer training is much more pronounced.

The best obtained errors are only slightly larger than in the previous case,
which indicates that the PINN still scales reasonably well up to this dimension.

For the two best-performing networks (80k):
\begin{center}
\begin{tabular}{ |c|c|c|c|c|c| } 
\hline
network size & $\ell_{\mathrm{total}}$ & rel $L^2$ & $L^\infty$ & peak allocated VRAM & runtime\\
\hline
4x256 & 1.126e-7 & 0.00215 & 0.00349 & 564 MB & 30 min\\
\hline
4x128 & 3.88e-8 & 0.00231 & 0.00376 & 464 MB & 27 min\\
\hline
\end{tabular}
\end{center}
% cluster_4d/

Before moving to even higher dimensions, we also test whether the training can be made more efficient.
In particular, we consider residual-based adaptive resampling (RAR) to focus on high-importance regions,
and stochastic-dimension gradient descent as a possible ways to reduce computational cost.

For these tests, we use the smaller 4x128 network trained for 80k steps and try to improve its performance without increasing its capacity.

First,
we consider stochastic-dimension gradient descent (SDGD) as a possible speedup.
Sampling only 2 dimensions gives at most a few percent runtime improvement, while noticeably worsening the loss landscape.
In this setup, this does not offer a beneficial tradeoff.

Next, we seek to optimize the training by focusing the network on the regions where it struggles.
In the RAR setting, we sample 4 times as many collocation points and then keep the top quarter with the highest residual according to the RAR formula.
We perform a parameter sweep over the corresponding values of $k$ and $c$.

Among the tested RAR settings, $k=2.0$ and $c=0.0$ gave the most promising behavior.

\begin{figure}
    \centering
    \IfFileExists{figures/nn_scaling_comparison_d=6.png}{
        \includegraphics[width=\linewidth]{figures/nn_scaling_comparison_d=6.png}
    }{
        \fbox{\parbox[c][0.3\textheight][c]{0.95\linewidth}{\centering Missing figure: \texttt{figures/nn\_scaling\_comparison\_d=6.png}}}
    }
    \caption{$d=6$: Error and loss for various hyperparameter choices of $k$ and $c$ in residual-based adaptive resampling (RAR).}
    \label{fig:models_d6}
\end{figure}

RAR has a relatively small effect on the model performance,


but it is biased towards minimizing the $L^\infty$ error.
which is sensible since it explicitly biases the point selection toward difficult regions.
The effect on the relative $L^2$ error is smaller, but still consistent with a modest gain in overall accuracy.



% mistake in running..., very good - same 1024 bs, new data every 100 steps
% what bout periodically giving the nn new data?


%\subsection{8th dimension}


\subsection{8 dimensions}

%If we continued with the patterns from the previous two dimentisions, we would next
%explore the networks of size 

As we saw in the previous two sections discussing $d=6$ and $d=4$, increasing the number of steps
improves the accuracy of the model.
However, such a change inevitably leads to a longer training run.

On the other hand, it is also apparent from Figures \ref{fig:models_d4} and \ref{fig:models_d6}
that, under a fixed number of steps, a larger network generally performs better,
while, crucially, finishing the simulation at comparable runtime.

In addition, in the most resource intensive run we have performed so far,
the VRAM peaked at only about 500 MB.

From this we conclude that the GPU is being underutilized.
Hence, there is still room to increase the network and batch size
while increasing the training time to a manageable degree.

In summary, 
we seek to map the problem of a long training run to one that we can address
by exploiting the advantages of the available hardware.

To this end, we consider two neural networks of size 4x512.
One with the batch size of 4096, the other of 8192.
Crucially, the training runs only for the same 80k steps as in the previous case, $d=6$.

\begin{center}
\begin{tabular}{ |c|c|c|c|c|c| } 
\hline
batch size & $\ell_{\mathrm{total}}$ & rel $L^2$ & $L^\infty$ & peak allocated VRAM & runtime\\
\hline
4096 & 1.565e-9 & 0.00188 & 0.0030 & 2.2 GB & 56 min\\
\hline
8192 & 2.392e-9 & 0.00198 & 0.0026 & 3.8 GB & 1h 40 min\\
\hline
\end{tabular}
\end{center}

Interestingly enough, if we had proceeded here as we did in the previous chapter, that is,
double the network width and the training domain, we would actually achive a singlificanly worse result, and a much slower one as well:
\begin{center}
\begin{tabular}{ |c|c|c|c|c|c| } 
\hline
batch size & $\ell_{\mathrm{total}}$ & rel $L^2$ & $L^\infty$ & peak allocated VRAM & runtime\\
\hline
4x256 & 4.110e-9 & 0.00440 & 0.01040 & 661 MB & 1h 48 min\\
\hline
4x512 & 4.274e-9 & 0.00521 & 0.01227 & 8.2 GB & 2h 44 min\\
\hline
\end{tabular}
\end{center}


At this point though, the training time still remains a burden.
When trianing PINNs, one usually has to do a few parameters sweeps to arrive at a satisfactory results,
and the cost of those sweeps rises with dimension.

%Training time remains the major bottleneck.
%At the same number of steps, a network with higher capacity performs better, as seen already in the scaling figures for $d=4$ and $d=6$.
%This suggests a natural next direction: using a larger network while reducing the total number of training steps.
%
%So far, the increase in required training steps is roughly linear to exponential,
%and the same is true for the model size.
%Scaling alone gets us surprisingly far, in a way quite similiar of scaling-law behavior,
%but the price is very long training.
%
%As a base configuration, we consider the 4x256 network and training horizons of 80k, 100k, and 160k steps.
%We then augment this setup with the lessons learned in 6 dimensions.
%
%At this point, however, the training cost becomes the dominant issue, and each additional gain in accuracy requires a visibly larger compute budget.


%As we can see from the previous 
%Exploit the advantages of GPU, parallel computation.
%Since we are still considerting relatively small networks and batch sizes,
%we can increse both and still achive a reasonable runtime.
%batches, we can increase the batch size and model capacity and
%relatively mild

%\begin{center}
%\begin{tabular}{ |c|c|c|c|c|c| } 
%\hline
%network size & $\ell_{\mathrm{total}}$ & rel $L^2$ & $L^\infty$ & peak allocated VRAM & runtime\\
%\hline
%4x256 & 4.110e-9 & 0.00440 & 0.01040 & 661 MB & 1h 48 min\\
%\hline
%4x512 & 4.274e-9 & 0.00521 & 0.01227 & 8.2 GB & 2h 44 min\\
%\hline
%\end{tabular}
%\end{center}


%\begin{center}
%\begin{tabular}{ |c|c|c|c|c|c| } 
%\hline
%network size & $\ell_{\mathrm{total}}$ & rel $L^2$ & $L^\infty$ & peak allocated VRAM & runtime\\
%\hline
%4x256 & 4.110e-9 & 0.00440 & 0.01040 & 661 MB & 1h 48 min\\
%\hline
%4x512 & 4.274e-9 & 0.00521 & 0.01227 & 8.2 GB & 2h 44 min\\
%\hline
%\end{tabular}
%\end{center}
%4096 & 1.565e-9 & 0.00188 & 0.0030 & 2.2 GB & 56 min\\
%8192 & 2.392e-9 & 0.00198 & 0.0026 & 3.8 GB & 1h 40 min\\






%\section{Curse of dimensionality}
%
%Both approaches clearly feel the curse of dimensionality, but in different ways.
%For sparse grids, the bottleneck is mainly combinatorial growth in the number of component grids and the associated memory cost.
%For PINNs, the bottleneck is less abrupt, but it appears through the steadily growing need for larger networks, longer training horizons, and more careful sampling strategies.
%
%The sparse-grid combination technique provides a strong method in low to moderate dimensions,
%but its computational cost rises quickly with the dimension and the level required to maintain an acceptable error tolerance.
%As the level or dimension increases, the memory needs of each worker rise as well.
%In higher dimensions, somewhat contrary to the parallel advantage that the combination technique offers,
%one must either dramatically increase the available memory or reduce the number of workers to fit the
%growing number of component grids and grid points in memory.
%
%
%It is also instructive to compare the degrees of freedom needed by the different approaches to achieve a given accuracy.
%In lower dimensions, sparse grids are highly competitive and often more direct to interpret.
%As the dimension grows, however, the PINN approach appears to degrade more gracefully, even though this comes at the cost of considerably longer and less predictable training.
%One important difference, however, is that the sparse-grid solver can be treated as a relatively direct numerical method once its discretization is fixed.
%By contrast, the PINN approach requires substantially more tuning, and therefore a fair comparison is not simply one sparse-grid run against one PINN run, but rather a sparse-grid run against a whole PINN model-selection process.
%
%
%\begin{figure}
%    \centering
%    \includegraphics[width=\linewidth]{figures/threshold_cost_comparison_rel_l2.png}
%    \caption{PINN vs combination technique perforamance in higher dimensions.
%    In the context of PINNs, the degrees of freedom denote the number of trainable parameters.
%    For the combination technique, it counts the total number of grid points in all component grids.}
%    \label{fig:comparison}
%\end{figure}




%\chapter{Appendix}
%% https://chatgpt.com/share/6a481b37-c04c-83eb-a0e8-b34c1d0ac08c
%
%\section{Heat equation: Derivation of the \texorpdfstring{$L^2$}{L2} norm of the analytic solution}
%Consider the analytical solution
%\begin{align*}
%u(t,x)
%&=
%e^{-\gamma t}
%\prod_{i=1}^d \sin(\pi x_i)
%+
%e^{-\delta t}
%\frac{0.67}{\sqrt d}
%\sum_{i=1}^d
%\sin(2\pi x_i)
%\left(
%\prod_{\substack{k=1 \\ k\neq i}}^d \sin(\pi x_k)
%\right)
%\end{align*}
%where
%\[
%\Omega = [0,1]^d
%\qquad
%\gamma = D\pi^2
%\qquad
%\delta = D\left(1+\frac 3d\right)\pi^2
%\]
%Here, we derive an explicit expression for the $L^2$ norm
%\[
%\|u(t,\cdot)\|_{L^2(\Omega)}^2
%=
%\int_{[0,1]^d} |u(t,x)|^2\,dx
%\]
%Define
%\[
%\phi_0(x)
%=
%\prod_{i=1}^d \sin(\pi x_i)
%\]
%and
%\[
%\phi_i(x)
%=
%\sin(2\pi x_i)
%\prod_{\substack{k=1 \\ k\neq i}}^d \sin(\pi x_k)
%\]
%Then the solution can be written compactly as
%\[
%u(t,x)
%=
%e^{-\gamma t}\phi_0(x)
%+
%\frac{0.67}{\sqrt d}e^{-\delta t}
%\sum_{i=1}^d \phi_i(x)
%\]
%Therefore
%\begin{align*}
%\|u(t,\cdot)\|_{L^2}^2
%&=
%\left\|
% e^{-\gamma t}\phi_0
% +
% \frac{0.67}{\sqrt d}e^{-\delta t}
% \sum_{i=1}^d \phi_i
%\right\|_{L^2}^2
%\end{align*}
%Expanding the expression gives
%\begin{align*}
%\|u(t,\cdot)\|_{L^2}^2
%&=
%e^{-2\gamma t}\|\phi_0\|_{L^2}^2
%+
%2e^{-\gamma t}e^{-\delta t}\frac{0.67}{\sqrt d}
%\sum_{i=1}^d \langle \phi_0,\phi_i\rangle_{L^2}
%\notag \\
%&\quad+
%\frac{0.67^2}{d}e^{-2\delta t}
%\sum_{i=1}^d\sum_{j=1}^d
%\langle \phi_i,\phi_j\rangle_{L^2}
%\end{align*}
%We now use the orthogonality relation
%\[
%\int_0^1 \sin(m\pi x)\sin(n\pi x)\,dx
%=
%\begin{cases}
%0, & m\neq n \\
%\dfrac12, & m=n
%\end{cases}
%\]
%Hence
%\begin{align*}
%\langle \phi_0,\phi_i\rangle_{L^2}&=0
%\qquad \text{for all } i=1,\dots,d \\
%\langle \phi_i,\phi_j\rangle_{L^2}&=0
%\qquad \text{for } i\neq j
%\end{align*}
%The remaining terms give
%\begin{equation*}
%\|\phi_0\|_{L^2}^2
%=
%\int_{[0,1]^d}
%\prod_{i=1}^d \sin^2(\pi x_i)\,dx
%=
%\prod_{i=1}^d
%\int_0^1 \sin^2(\pi x_i)\,dx_i
%=
%\frac{1}{2^d}
%\end{equation*}
%
%\begin{equation*}
%\|\phi_i\|_{L^2}^2
%=
%\left(
%\int_0^1 \sin^2(2\pi x_i)\,dx_i
%\right)
%\prod_{\substack{k=1 \\ k\neq i}}^d
%\left(
%\int_0^1 \sin^2(\pi x_k)\,dx_k
%\right)
%=
%\frac{1}{2^d}
%\end{equation*}
%
%\begin{equation*}
%\sum_{i=1}^d\sum_{j=1}^d
%\langle \phi_i,\phi_j\rangle_{L^2}
%=
%\sum_{i=1}^d \|\phi_i\|_{L^2}^2
%=
%d\frac{1}{2^d}
%\end{equation*}
%Substituting into the orignal expression, we obtain
%\begin{align*}
%\|u(t,\cdot)\|_{L^2}^2
%&=
%e^{-2\gamma t}\frac{1}{2^d}
%+
%\frac{0.67^2}{d}e^{-2\delta t}
%\,
%d\frac{1}{2^d}
%\end{align*}
%Thanks to the normalization $1/\sqrt d$, the factor $d$ cancels.
%Therefore
%\[
%\|u(t,\cdot)\|_{L^2}^2
%=
%\frac{1}{2^d}
%\left(
% e^{-2\gamma t}
% +
% 0.67^2 e^{-2\delta t}
%\right)
%\]
%At $t=0$, this simplifies to
%\[
%\|u_0\|_{L^2}^2
%=
%\frac{1}{2^d}
%(1+0.67^2)
%\]
%
